% GENERATED FILE. Edit canonical lessons, then run npm run generate:books.
% canonical-source-hash: fnv1a64:81cbee930832f5fd
% canonical-lessons: RU-C114-price, RU-C114-football, RU-C114-concert, RU-C114-holiday, RU-C114-fog

\chapter{Price, Football, Concert, Holiday, Fog}
\label{ch:ru-price-football-concert-holiday-fog}
\hlchaptermodality{114}

\begin{chapteropening}
\textbf{By the end of this chapter:} \emph{I can say a price, football, a concert, a holiday and fog.}

Complete the last lesson of chapter 114: I can say a price, football, a concert, a holiday and fog.
\end{chapteropening}

\section[\texorpdfstring{\ru{цена} (tsená)}{цена (tsená)}]{\ru{цена} (tsená) — a price}

\label{lesson:RU-C114-price}



\begin{quote}
Before the new one: say the Russian for a drawing, then the Russian for life.
\end{quote}

\subsection*{You'll want to know: \ru{цена}}
\textbf{\ru{цена}} (\emph{tsená}) — \textquotedblleft{}a price\textquotedblright{}.

\begin{cousinweb}
\textbf{\ru{цена}} $\leftarrow$ Proto-Slavic \emph{*cěna} — inherited. Vasmer relates it to Lithuanian \emph{kainà}, \textquotedblleft{}price\textquotedblright{}, Avestan \emph{kaēnā-}, \textquotedblleft{}retribution\textquotedblright{}, and Greek \emph{poin\'{\={e}}}, \textquotedblleft{}penalty\textquotedblright{}, the source of English \emph{pain}.
\end{cousinweb}

\begin{grammarlens}[title={What you've built}]
One more word for this chapter.
\end{grammarlens}

\subsection*{Your turn}
\begin{itemize}
\raggedright
  \item \emph{Say it:} \emph{tsená}
  \item \emph{Say it:} \emph{tsená}, once more
  \item \emph{Say it:} say \emph{tsená}
  \item \emph{Recall:} say the Russian for a drawing, then the Russian for life, then say \emph{tsená} again
\end{itemize}

\begin{tcolorbox}[breakable,colback=teal!4,colframe=teal!35!black,title={Before you move on}]
What does \ru{цена} mean? (\textquotedblleft{}A price\textquotedblright{}.) Say it once more.
\end{tcolorbox}
\section[\texorpdfstring{\ru{футбол} (futból)}{футбол (futból)}]{\ru{футбол} (futból) — football}

\label{lesson:RU-C114-football}



\begin{quote}
Before the new one: say the Russian for life, then the Russian for a price.
\end{quote}

\subsection*{You'll want to know: \ru{футбол}}
\textbf{\ru{футбол}} (\emph{futból}) — \textquotedblleft{}football\textquotedblright{}.

\begin{cousinweb}
\textbf{\ru{футбол}} comes from English \emph{football}, \textquotedblleft{}foot ball\textquotedblright{} (Trubachev, in Vasmer).
\end{cousinweb}

\begin{grammarlens}[title={What you've built}]
One more word for this chapter.
\end{grammarlens}

\subsection*{Your turn}
\begin{itemize}
\raggedright
  \item \emph{Say it:} \emph{futból}
  \item \emph{Say it:} \emph{futból}, once more
  \item \emph{Say it:} say \emph{futból}
  \item \emph{Recall:} say the Russian for life, then the Russian for a price, then say \emph{futból} again
\end{itemize}

\begin{tcolorbox}[breakable,colback=teal!4,colframe=teal!35!black,title={Before you move on}]
What does \ru{футбол} mean? (\textquotedblleft{}Football\textquotedblright{}.) Say it once more.
\end{tcolorbox}
\section[\texorpdfstring{\ru{концерт} (kontsért)}{концерт (kontsért)}]{\ru{концерт} (kontsért) — a concert}

\label{lesson:RU-C114-concert}



\begin{quote}
Before the new one: say the Russian for a price, then the Russian for football.
\end{quote}

\subsection*{You'll want to know: \ru{концерт}}
\textbf{\ru{концерт}} (\emph{kontsért}) — \textquotedblleft{}a concert\textquotedblright{}.

\begin{cousinweb}
\textbf{\ru{концерт}} came through German \emph{Konzert} from Italian \emph{concerto}, \textquotedblleft{}a contest (of voices)\textquotedblright{}, from Latin \emph{concertāre}, \textquotedblleft{}to contend\textquotedblright{}; under Peter the Great it meant \textquotedblleft{}an agreement\textquotedblright{} (Vasmer).
\end{cousinweb}

\begin{grammarlens}[title={What you've built}]
One more word for this chapter.
\end{grammarlens}

\subsection*{Your turn}
\begin{itemize}
\raggedright
  \item \emph{Say it:} \emph{kontsért}
  \item \emph{Say it:} \emph{kontsért}, once more
  \item \emph{Say it:} say \emph{kontsért}
  \item \emph{Recall:} say the Russian for a price, then the Russian for football, then say \emph{kontsért} again
\end{itemize}

\begin{tcolorbox}[breakable,colback=teal!4,colframe=teal!35!black,title={Before you move on}]
What does \ru{концерт} mean? (\textquotedblleft{}A concert\textquotedblright{}.) Say it once more.
\end{tcolorbox}
\section[\texorpdfstring{\ru{праздник} (prázdnik)}{праздник (prázdnik)}]{\ru{праздник} (prázdnik) — a holiday, a celebration}

\label{lesson:RU-C114-holiday}



\begin{quote}
Before the new one: say the Russian for football, then the Russian for a concert.
\end{quote}

\subsection*{You'll want to know: \ru{праздник}}
\textbf{\ru{праздник}} (\emph{prázdnik}) — \textquotedblleft{}a holiday, a celebration\textquotedblright{}.

\begin{cousinweb}
\textbf{\ru{праздник}} is built on \emph{\ru{праздный}}, \textquotedblleft{}idle, free\textquotedblright{}, which Vasmer calls Church Slavonic, by its \emph{-\ru{ра}-}, beside native \emph{\ru{порожний}}, \textquotedblleft{}empty\textquotedblright{}: a holiday is a day free of work.
\end{cousinweb}

\begin{grammarlens}[title={What you've built}]
One more word for this chapter.
\end{grammarlens}

\subsection*{Your turn}
\begin{itemize}
\raggedright
  \item \emph{Say it:} \emph{prázdnik}
  \item \emph{Say it:} \emph{prázdnik}, once more
  \item \emph{Say it:} say \emph{prázdnik}
  \item \emph{Recall:} say the Russian for football, then the Russian for a concert, then say \emph{prázdnik} again
\end{itemize}

\begin{tcolorbox}[breakable,colback=teal!4,colframe=teal!35!black,title={Before you move on}]
What does \ru{праздник} mean? (\textquotedblleft{}A holiday, a celebration\textquotedblright{}.) Say it once more.
\end{tcolorbox}
\section[\texorpdfstring{\ru{туман} (tumán)}{туман (tumán)}]{\ru{туман} (tumán) — fog}

\label{lesson:RU-C114-fog}



\begin{quote}
Before the new one: say the Russian for a concert, then the Russian for a holiday.
\end{quote}

\subsection*{You'll want to know: \ru{туман}}
\textbf{\ru{туман}} (\emph{tumán}) — \textquotedblleft{}fog\textquotedblright{}.

\begin{cousinweb}
\textbf{\ru{туман}} comes from Turkic \emph{tuman}, \textquotedblleft{}fog, gloom\textquotedblright{}, as in Kazakh and Kirghiz, Turkish \emph{duman}; Vasmer thinks the Turkic words probably go back to Avestan \emph{dunman-}, \textquotedblleft{}mist\textquotedblright{}.
\end{cousinweb}

\begin{grammarlens}[title={What you've built}]
One more word for this chapter.
\end{grammarlens}

\subsection*{Your turn}
\begin{itemize}
\raggedright
  \item \emph{Say it:} \emph{tumán}
  \item \emph{Say it:} \emph{tumán}, once more
  \item \emph{Say it:} say \emph{tumán}
  \item \emph{Recall:} say the Russian for a concert, then the Russian for a holiday, then say \emph{tumán} again
\end{itemize}

\begin{tcolorbox}[breakable,colback=teal!4,colframe=teal!35!black,title={Before you move on}]
What does \ru{туман} mean? (\textquotedblleft{}Fog\textquotedblright{}.) Say it once more.
\end{tcolorbox}
